% code-contact-map-generator.tex
% Code Contact Map Generator: A Measure-Theoretic Framework for Minimum
% Sufficient Mutual Individuation in Software Systems
%
% Compile with: pdflatex -> bibtex -> pdflatex -> pdflatex

\documentclass[12pt,a4paper,twocolumn]{article}

% --------------------------------------------------------------------------
% Packages
% --------------------------------------------------------------------------
\usepackage[a4paper, margin=2cm, columnsep=0.8cm]{geometry}
\usepackage{amsmath,amssymb,amsthm}
\usepackage{mathtools}
\usepackage{bbm}
\usepackage{stmaryrd}        % \llbracket, \rrbracket
\usepackage{graphicx}
\usepackage{xcolor}
\usepackage{hyperref}
\usepackage{natbib}
\usepackage{listings}
\usepackage{microtype}
\usepackage{booktabs}
\usepackage{enumitem}
\usepackage{float}
\usepackage{caption}
\usepackage{subcaption}

% --------------------------------------------------------------------------
% Hyperref setup
% --------------------------------------------------------------------------
\hypersetup{
  colorlinks=true,
  linkcolor=blue!60!black,
  citecolor=green!50!black,
  urlcolor=blue!70!black,
  pdftitle={Code Contact Map Generator},
  pdfauthor={},
}

% --------------------------------------------------------------------------
% Listings / DSL syntax
% --------------------------------------------------------------------------
\definecolor{dsl-keyword}{RGB}{0,100,200}
\definecolor{dsl-string}{RGB}{160,40,40}
\definecolor{dsl-comment}{RGB}{80,120,80}
\definecolor{dsl-number}{RGB}{140,60,0}
\definecolor{dsl-bg}{RGB}{248,248,252}
\definecolor{dsl-frame}{RGB}{200,200,220}

\lstdefinelanguage{WindTunnel}{
  morekeywords={scope,analyse,compose,assert,recurse,report,
                target,exclude,method,model,query,merge,resolve,
                extract,claim,on_fail,output,format,until,while,
                static,dynamic,semantic,hf,individuation,regime,
                distance,narrow_scope,marginal_return,floor,by},
  sensitive=true,
  morestring=[b]",
  morestring=[b]',
  morecomment=[l]{\#},
}

\lstset{
  language=WindTunnel,
  backgroundcolor=\color{dsl-bg},
  frame=single,
  framerule=0.4pt,
  rulecolor=\color{dsl-frame},
  basicstyle=\ttfamily\footnotesize,
  keywordstyle=\bfseries\color{dsl-keyword},
  stringstyle=\color{dsl-string},
  commentstyle=\itshape\color{dsl-comment},
  numberstyle=\tiny\color{gray},
  numbers=left,
  numbersep=5pt,
  breaklines=true,
  breakatwhitespace=false,
  captionpos=b,
  tabsize=2,
  showspaces=false,
  showstringspaces=false,
}

% --------------------------------------------------------------------------
% Theorem environments
% --------------------------------------------------------------------------
\theoremstyle{plain}
\newtheorem{theorem}{Theorem}[section]
\newtheorem{lemma}[theorem]{Lemma}
\newtheorem{proposition}[theorem]{Proposition}
\newtheorem{corollary}[theorem]{Corollary}

\theoremstyle{definition}
\newtheorem{definition}[theorem]{Definition}
\newtheorem{example}[theorem]{Example}

\theoremstyle{remark}
\newtheorem{remark}[theorem]{Remark}

% --------------------------------------------------------------------------
% Notation macros
% --------------------------------------------------------------------------
\newcommand{\Borel}{\mathcal{B}}
\newcommand{\Part}{\mathcal{P}}
\newcommand{\Sep}{\mathrm{Sep}}
\newcommand{\CM}{\mathrm{CM}}
\newcommand{\bmin}{\beta_{\min}}
\newcommand{\bP}{\beta_{\mathcal{P}}}
\newcommand{\Scoord}{S}
\newcommand{\dS}{d_S}
\newcommand{\Neg}{\mathrm{Neg}}
\newcommand{\Rest}{\mathrm{R}_{\mathrm{est}}}
\newcommand{\Lmin}{L_{\min}}
\newcommand{\Smin}{S_{\min}}
\newcommand{\CMbmin}{\CM_{\bmin}}
\newcommand{\GC}{G_C}
\newcommand{\GS}{G_S}
\newcommand{\hol}{\mathrm{hol}}
\newcommand{\diag}{\mathrm{diag}}
\newcommand{\val}{\mathrm{val}}
\newcommand{\Sbmin}{S_{\bmin}}
\newcommand{\Real}{\mathbb{R}}
\newcommand{\Nat}{\mathbb{N}}
\newcommand{\BRS}{(\Omega,d,\mu)}

% --------------------------------------------------------------------------
% Title
% --------------------------------------------------------------------------
\title{\textbf{Code Contact Map Generator}\\[4pt]
  {\large A Measure-Theoretic Framework for Minimum Sufficient\\
  Mutual Individuation in Software Systems}}

\author{Wind Tunnel Research\\[4pt]
  \small \texttt{wind-tunnel.dev}}

\date{June 2026}

% ==========================================================================
\begin{document}
% ==========================================================================

\maketitle

% --------------------------------------------------------------------------
\begin{abstract}
% --------------------------------------------------------------------------
We introduce the \emph{Code Contact Map Generator}, a framework for
characterising software systems not by the binary outcomes of individual test
assertions, but by the \emph{minimum sufficient contact map} --- the minimal
measure-theoretic structure that certifies mutual individuation across every
pair of system components.  The foundational object is the \emph{bounded
resolvable space} $\BRS$, a compact metric space equipped with a finite
non-atomic Borel measure and a resolution floor.  Within this space every
valid partition of the system into components must deposit a non-zero
separator between each pair of adjacent components.  We derive the
\emph{partition floor} $\bmin > 0$ from first principles: an agent is an
individuated item within the space; individuation requires comparison against
an infinite complement; each comparison step must deposit a cell of
positive measure; therefore $\bmin > 0$ for any finite agent operating in
an infinite space.  This derivation replaces earlier geometric proofs and
grounds the floor in the structure of individuation itself --- not in
computational budget limits.  We derive from this the \emph{Minimum
Sufficiency Trijection}: the partition floor $\bmin$, the minimum sufficient
contact map $\CMbmin$, and the minimum sufficient Lagrangian $\Lmin$ are in
canonical bijective correspondence; the minimum sufficient action $\Smin$ is
a derived quantity (the integral of $\Lmin$), not a fourth independent
object.  We present a new section establishing the empirical grounding of
the framework through the structure of scientific measurement: truth is a
cell, not a point, and every instrument array in science is a contact map
approximation from below.  We then model software systems as bounded
resolvable spaces, define S-entropy coordinates for code units, and use them
to derive five coordination regimes (operationally calibrated) quantified by
the static order parameter $\Rest$.  The framework is realised by the
\emph{Wind Tunnel DSL}, a partition-specification language in which both
users and the contact-map generator co-author investigation scripts.
\end{abstract}

\tableofcontents

\onecolumn   % switch to single column for the body

% ==========================================================================
\section{Introduction}
% ==========================================================================

Software verification has, since its inception, been organised around a
single epistemic act: the evaluation of a predicate at a point.  A test is
executed, a return value is inspected, a property is declared true or false.
This paradigm is so deeply embedded in engineering practice that it shapes not
only tooling but conceptual vocabulary: ``passing the tests,'' ``green build,''
``coverage report.''  The implicit model is that correctness is a point in a
binary space, and that enough point evaluations constitute a proof.

This paper challenges that model at its foundation.  We argue that the
natural unit of software-system characterisation is not the truth value of an
isolated assertion but the \emph{contact map}: the assignment, to each pair of
system components, of the minimum information required to distinguish them from
each other.  A test failure is evidence that contributes to the contact map;
it is not the map itself.  The contact map is a global structure; point
evaluations are local observations.  The relationship is analogous to the
distinction between a measure and the value of its density at a single point.

\subsection{The Problem with Truth-at-a-Point}

Let $f : X \to Y$ be a software function and $\phi : X \times Y \to
\{0,1\}$ a specification predicate.  Classical testing evaluates $\phi(x_i,
f(x_i))$ for a finite sample $\{x_i\}_{i=1}^N \subset X$.  Several
fundamental obstacles limit the conclusions available from this procedure.

\begin{enumerate}[label=(\roman*)]
\item \textbf{Rice's theorem} \citep{rice1953classes}: no non-trivial semantic
  property of a program is decidable.  Any complete automated oracle is
  impossible.

\item \textbf{Undecidability of halting} \citep{turing1936computable}: the
  test harness itself cannot determine, in general, whether the system under
  test terminates.

\item \textbf{Local/global gap}: $\phi(x_i, f(x_i)) = 1$ for all sampled
  $x_i$ does not imply $\forall x \in X\!: \phi(x,f(x)) = 1$.  This is not
  merely a practical limitation; it is a categorical one.

\item \textbf{Compositionality failure}: local correctness of each component
  does not imply global correctness of the composed system (formalised in
  Theorem~\ref{thm:local-invisibility} below).
\end{enumerate}

These obstacles share a common source: the point-evaluation paradigm attempts
to characterise a measure-theoretic object (the behaviour of $f$ over a
continuous or large discrete domain) by a finite sequence of point samples.
The information deposited by each sample is bounded; the object is not.

\subsection{Contact Map Generation as an Alternative}

We propose replacing the question ``did this unit pass or fail?'' with the
question ``what is the minimum sufficient structure of mutual individuation
across this system?''  These are categorically different questions.  The
first is local and binary; the second is global and metric.  The framework
developed in this paper provides:

\begin{enumerate}[label=(\roman*)]
\item A \textbf{mathematical foundation} (Sections~\ref{sec:foundations}
  and~\ref{sec:mscm}) based on measure theory and compact metric spaces,
  yielding the \emph{partition floor} $\bmin > 0$ derived from the
  structure of individuation in an infinite space.

\item An \textbf{empirical grounding} (Section~\ref{sec:empirical}) showing
  that every instrument array in science is a contact map approximation from
  below, and that truth is structurally a cell, not a point.

\item A \textbf{model of software systems} (Section~\ref{sec:software}) as
  bounded resolvable spaces, with S-entropy coordinates for code units and
  five quantified coordination regimes.

\item A \textbf{contact map generator} (Section~\ref{sec:generator}) that
  uses static analysis, dynamic tracing, and semantic inference as partition
  operations contributing estimates of $\CMbmin$.

\item The \textbf{Wind Tunnel DSL} (Section~\ref{sec:dsl}), a partition-
  specification language with formal grammar and compositional operational
  semantics, in which both users and the system co-author investigation
  scripts.
\end{enumerate}

\subsection{Overview and Contributions}

The central theoretical result is the \emph{Minimum Sufficiency Trijection}
(Theorem~\ref{thm:minimum-sufficiency}), which establishes a canonical
bijective correspondence between three objects: the partition floor, the
minimum sufficient contact map, and the minimum sufficient Lagrangian.  The
minimum sufficient action is a derived corollary.  All additional structural
annotations are \emph{partition depth inflations} of this ground state.

The central methodological result is the \emph{Local Invisibility Theorem}
(Theorem~\ref{thm:local-invisibility}): universal local validity does not
imply global contact map coherence.  This is the formal basis for why
contact map generation is strictly more informative than test aggregation.

The \emph{Wind Tunnel DSL} (Section~\ref{sec:dsl}) is the language in which
both questions --- classical test assertions and contact map claims --- can be
expressed.  A DSL execution file is a complete, reproducible record of a
contact map investigation, suitable for auditing, extension, and collaborative
authorship between human and automated agents.

% ==========================================================================
\section{Mathematical Foundations}
\label{sec:foundations}
% ==========================================================================

\subsection{Bounded Resolvable Spaces}

\begin{definition}[Bounded Resolvable Space]
\label{def:brs}
A \emph{bounded resolvable space} is a triple $\BRS$ where:
\begin{enumerate}[label=(\roman*)]
  \item $(\Omega, d)$ is a compact metric space;
  \item $\mu$ is a finite non-atomic Borel measure on $\Omega$ with
    $\mu(\Omega) < \infty$;
  \item There exist constants $\delta, c_0 > 0$ and $n \geq 1$ such that
    for every $x \in \Omega$ and every $0 < r < \delta$,
    \begin{equation}
      \mu\!\left(B(x,r)\right) \;\geq\; c_0\, r^n,
      \label{eq:resolution-floor}
    \end{equation}
    where $B(x,r) = \{y \in \Omega : d(x,y) < r\}$ is the open $d$-ball.
\end{enumerate}
The triple $(\delta, c_0, n)$ is called the \emph{resolution data}, and
condition~\eqref{eq:resolution-floor} is the \emph{resolution floor}.
\end{definition}

\begin{remark}
The resolution floor is a lower Ahlfors regularity condition
\citep{mattila1995geometry, falconer2003fractal}: every ball of radius $r$
carries at least $c_0 r^n$ measure.  It excludes purely singular or purely
atomic measures and ensures that every open set of positive diameter carries
positive measure.  Euclidean space with Lebesgue measure satisfies this with
$n = \dim$; self-similar fractals satisfying the open set condition satisfy
it with $n$ equal to the Hausdorff dimension.
\end{remark}

\subsection{Partitions and Partition Depth}

\begin{definition}[Valid Partition]
\label{def:partition}
A \emph{valid partition} of $\BRS$ is a finite collection
$\Part = \{A_1, \ldots, A_k\}$ of measurable subsets of $\Omega$ satisfying:
\begin{enumerate}[label=(\roman*)]
  \item \textbf{Covering}: $\mu\!\left(\Omega \setminus \bigcup_{i=1}^k A_i\right) = 0$;
  \item \textbf{Essential disjointness}: $\mu(A_i \cap A_j) = 0$ for $i \neq j$;
  \item \textbf{Connectedness}: each $A_i$ is connected in the subspace
    topology induced by $d$;
  \item \textbf{Positive measure}: $\mu(A_i) > 0$ for all $i$.
\end{enumerate}
Denote the set of all valid partitions by $\mathfrak{P}(\Omega,d,\mu)$.
\end{definition}

\begin{definition}[Separator]
\label{def:separator}
For $A \in \Part$ with $\Part \in \mathfrak{P}$, the \emph{separator set}
of $A$ in $\Part$ is
\begin{equation}
  \Sep(A, \Part) \;=\; \bigl\{ B \in \Part \setminus \{A\}
    : \overline{A} \cap \overline{B} \neq \emptyset \bigr\},
\end{equation}
the collection of components whose topological closures share a boundary
point with $\overline{A}$.
\end{definition}

\begin{definition}[Partition Depth]
\label{def:partition-depth}
The \emph{partition depth} of $A \in \Part$ is
\begin{equation}
  \bP(A) \;=\; \inf_{B \in \Sep(A,\Part)} \mu(B).
\end{equation}
The \emph{partition floor} of the space is
\begin{equation}
  \bmin \;=\; \inf_{\Part \in \mathfrak{P}}\; \inf_{A \in \Part}\; \bP(A).
\end{equation}
\end{definition}

\begin{remark}
The partition depth $\bP(A)$ is the smallest measure carried by any component
adjacent to $A$.  The partition floor $\bmin$ is the infimum of this quantity
over all possible partitions: the smallest separator any partition strategy
can produce.
\end{remark}

\subsection{Agents as Individuated Items}

Before proving the positivity of $\bmin$, we establish the correct setting
for the floor.  The notion of an ``agent'' --- an entity that performs
partition operations --- is not an external assumption about the space.  It
follows from the structure of individuation itself.

\begin{definition}[Agent]
\label{def:agent}
An \emph{agent} is an individuated item $\alpha \in \Part$ within a bounded
resolvable space $\BRS$ that performs further individuation operations on
other components of $\Omega$.  An agent is individuated by negation:
$\alpha$ is fixed as the unique $\mu$-a.e.\ element satisfying
$\alpha \cup \Neg(\alpha, \Part) = \Omega$ and $\mu(\alpha \cap
\Neg(\alpha, \Part)) = 0$.
\end{definition}

\begin{proposition}[Agents are Finite]
\label{prop:agents-finite}
Any agent $\alpha \in \Part$ within $\BRS$ satisfies $\mu(\alpha) < \mu(\Omega)$.
In particular, $\mu(\Neg(\alpha, \Part)) > 0$.
\end{proposition}

\begin{proof}
Suppose $\mu(\alpha) = \mu(\Omega)$.  Then $\mu(\Neg(\alpha, \Part)) =
\mu(\Omega \setminus \alpha) = 0$, which means $\Neg(\alpha, \Part)$ has
zero measure.  But the separator of $\alpha$ in any partition must include
at least one component of positive measure (by Definition~\ref{def:partition}(iv));
the separator is a subset of $\Neg(\alpha, \Part)$; hence the separator has
measure zero, violating Definition~\ref{def:partition-depth}.  Therefore
$\mu(\alpha) < \mu(\Omega)$, and $\mu(\Neg(\alpha, \Part)) > 0$.
\end{proof}

\begin{remark}
This proposition establishes that finiteness (in the sense of strict
sub-totality) is not an assumption about agents --- it is a consequence of
individuation.  An item that occupied the entire space would be
indistinguishable from the space itself: its negation would be empty,
violating the positive-separator condition.  Individuation by negation
enforces finiteness.
\end{remark}

\subsection{Individuation by Negation}

\begin{definition}[Negation of a Component]
\label{def:negation}
For $A \in \Part$, the \emph{negation} of $A$ in $\Part$ is
\begin{equation}
  \Neg(A, \Part) \;=\; \Omega \setminus A \;=\; \bigcup_{B \in \Part, B \neq A} B.
\end{equation}
\end{definition}

\begin{proposition}[Selector-Free Individuation]
\label{prop:selector-free}
The negation $\Neg(A,\Part)$ is the unique individuating complement of $A$
that requires no distinguished element or selector function.  Specifically,
$A$ is the $\mu$-a.e.\ unique set satisfying $A \cup \Neg(A,\Part) = \Omega$
and $\mu(A \cap \Neg(A,\Part)) = 0$.
\end{proposition}

\begin{proof}
By essential disjointness (Definition~\ref{def:partition}(ii)) and the
covering condition (Definition~\ref{def:partition}(i)), the complement
$\Omega \setminus A$ is well-defined up to a $\mu$-null set.  No choice
function is needed: the complement is simply the union of all other cells.
Uniqueness follows from the atomless property of $\mu$: if $A'$ also satisfies
$A' \cup (\Omega \setminus A) = \Omega$ and $\mu(A' \cap (\Omega \setminus
A)) = 0$, then $\mu(A \triangle A') = 0$, so $A = A'$ $\mu$-a.e.
\end{proof}

\begin{remark}
This formalises the core principle: a component $A$ is fixed as ``the whole
minus everything it is not.''  The complement carries no distinguished
element; it is the structural residue of $A$'s existence within $\Part$.
Individuation requires no internal reference point --- only the partition
and the complement.
\end{remark}

\subsection{Incompletability of Individuation}

The key structural fact about individuation in an infinite space is that the
comparison against the complement is never completable by a finite agent.

\begin{theorem}[Incompletability of Individuation]
\label{thm:incompletable}
Let $\BRS$ be a bounded resolvable space with $\mu$ non-atomic and
$\Omega$ infinite.  For any agent $\alpha \in \Part$ and any finite
sequence of partition operations $(\mathcal{O}^{(1)}, \ldots, \mathcal{O}^{(n)})$
performed by $\alpha$, the comparison of $\alpha$ against
$\Neg(\alpha, \Part)$ remains incomplete: there exists a measurable set
$R \subseteq \Neg(\alpha, \Part)$ with $\mu(R) > 0$ that has not been
individuated from $\alpha$ by the sequence.
\end{theorem}

\begin{proof}
Each partition operation $\mathcal{O}^{(j)}$ carves a cell $C_j$ from
$\Neg(\alpha, \Part)$ and certifies ``$\alpha$ is not $C_j$.''  By
Definition~\ref{def:partition}(iv), each $C_j$ has $\mu(C_j) > 0$.  After
$n$ operations, the certified region is $\bigcup_{j=1}^n C_j$ with
$\mu\!\bigl(\bigcup_j C_j\bigr) \leq \sum_j \mu(C_j) < \infty$.  Since
$\mu(\Neg(\alpha, \Part)) > 0$ (by Proposition~\ref{prop:agents-finite})
and $\mu$ is non-atomic on an infinite space, there exist uncountably many
measurable subsets of $\Neg(\alpha, \Part)$ with positive measure.  A
finite sequence of operations covers only finitely many.  Hence the
residual $R = \Neg(\alpha, \Part) \setminus \bigcup_j C_j$ satisfies
$\mu(R) > 0$.
\end{proof}

\begin{remark}
The incompletability is not a computational limitation; it is structural.
The complement $\Neg(\alpha, \Part)$ is ``as large as the universe minus
$\alpha$.''  Every finite sequence of partition operations certifies a
finite region of this complement.  The comparison terminates only when the
agent is no longer finite --- i.e., never.  This gives individuation a
strict direction: each step advances the committed record, but no finite
sequence of steps completes it.  The process is monotone and non-terminating:
exactly the conditions of the Non-Return Theorem below.
\end{remark}

\subsection{Positivity of the Partition Floor}

We now derive $\bmin > 0$ from the structure of individuation.  This
replaces earlier geometric or representational arguments, which required
unjustified inequalities about cell-size distribution.

\begin{theorem}[Partition Floor is Strictly Positive]
\label{thm:floor-positive}
Let $\BRS$ be a bounded resolvable space.  Then $\bmin > 0$.
\end{theorem}

\begin{proof}
Let $A \in \Part$ be any component in any valid partition
$\Part \in \mathfrak{P}$.  To establish $A$ as individual --- to certify
that $A$ is not some $B \in \Sep(A, \Part)$ --- requires performing a
partition operation that carves $B$ from $\Neg(A, \Part)$ and deposits $B$
into the separator of $A$.

A partition operation deposits a cell $C \subseteq \Neg(A, \Part)$ with
$\mu(C) > 0$ (by Definition~\ref{def:partition}(iv): cells have positive
measure; a cell of zero measure is not a cell --- it is engulfed by its own
boundary and undetectable).

The separator $\Sep(A, \Part)$ is the collection of such cells adjacent to
$A$.  Every element $B \in \Sep(A, \Part)$ is a cell, hence $\mu(B) > 0$.
Therefore:
\begin{equation}
  \bP(A) \;=\; \inf_{B \in \Sep(A,\Part)} \mu(B) \;\geq\;
    \min_{B \in \Sep(A,\Part)} \mu(B) \;>\; 0,
\end{equation}
where the last inequality holds because $|\Sep(A, \Part)| \leq |\Part| - 1 < \infty$
(the partition is finite by Definition~\ref{def:partition}) and each
$\mu(B) > 0$, so the minimum over a finite set of positive quantities is
positive.

Since this holds for every $A \in \Part$ and every $\Part \in \mathfrak{P}$,
we have $\bmin = \inf_{\Part} \inf_A \bP(A) > 0$ as required.
\end{proof}

\begin{remark}
The proof depends on two facts that are jointly equivalent to individuation
by negation in a bounded resolvable space: (1) cells must have positive
measure (the partition axioms); (2) the partition is finite (the agent is
finite, by Proposition~\ref{prop:agents-finite}).  Both are consequences of
the bounded resolvable space structure, not additional assumptions.  The
floor $\bmin$ is not a property of the investigation method; it is a
property of what it means to carve a cell from a complement.
\end{remark}

\begin{remark}[On the intrinsic vs.\ agent-relative floor]
The floor $\bmin$ as defined here is relative to finite valid partitions ---
that is, to the partitions realisable by any finite agent within the space.
This is the correct scope.  By Proposition~\ref{prop:agents-finite}, any
entity that could measure the floor is itself individuated, hence finite,
hence agent-relative.  There is no ``intrinsic floor'' accessible to an
infinite observer, because such an observer cannot be individuated within
the space.  The agent-relative floor is the only floor.
\end{remark}

\subsection{Detectability}

\begin{definition}[Signal Measure]
\label{def:signal}
For a component $A \in \Part$, its \emph{signal measure} is
\begin{equation}
  \sigma(A) \;=\; \mu(A) - \mu(\partial A),
\end{equation}
where $\partial A = \overline{A} \setminus \mathrm{int}(A)$ is the topological
boundary of $A$ in $(\Omega, d)$.
\end{definition}

\begin{definition}[Detectability]
\label{def:detectability}
A component $A \in \Part$ is \emph{detectable} if
\begin{equation}
  \sigma(A) \;>\; \bP(A).
\end{equation}
It is \emph{subdetectable} (engulfed by its own boundary) if
$\sigma(A) \leq \bP(A)$.
\end{definition}

\begin{proposition}
In a bounded resolvable space, a component $A$ is detectable if and only if
it can be distinguished, with positive $\mu$-probability, from any element
of $\Sep(A,\Part)$ by a measure-preserving observation process.
\end{proposition}

\begin{proof}
If $\sigma(A) > \bP(A)$, then $A$ carries strictly more interior measure
than the measure of any adjacent separator.  A measure-preserving process
visiting $A$ will encounter its interior --- not merely its boundary --- with
positive probability, enabling distinction.  Conversely, if $\sigma(A) \leq
\bP(A)$, then the boundary accounts for at least as much measure as any
adjacent component; observations of $A$ are indistinguishable from
observations of $\partial A$, which is shared with adjacent components.
\end{proof}

\subsection{Locating vs.\ Visiting Asymmetry}

\begin{definition}[Measure-Preserving Traversal]
\label{def:traversal}
A \emph{measure-preserving traversal} of $\BRS$ is a measure-preserving
ergodic transformation $T : \Omega \to \Omega$.  For a component $A$,
the \emph{visiting frequency} of $A$ under $T$ is
\begin{equation}
  \nu_T(A) \;=\; \lim_{N\to\infty} \frac{1}{N} \sum_{j=0}^{N-1}
    \mathbf{1}_A\!\left(T^j(x)\right) \;=\; \mu(A),
\end{equation}
where the last equality holds $\mu$-a.e.\ by the Birkhoff ergodic theorem
\citep{halmos1950measure}.
\end{definition}

\begin{definition}[Locating Cost]
\label{def:locating-cost}
The \emph{locating cost} of a component $A$ with visiting probability
$p_A = \mu(A)/\mu(\Omega)$ is
\begin{equation}
  \lambda(A) \;=\; \bmin \cdot \left\lceil \log_2 \frac{1}{p_A} \right\rceil.
\end{equation}
\end{definition}

\begin{theorem}[Locating/Visiting Asymmetry]
\label{thm:locating-visiting}
Let $T$ be a measure-preserving traversal of $\BRS$ and $A \in \Part$ a
component with $p_A = \mu(A)/\mu(\Omega)$.  Then:
\begin{enumerate}[label=(\roman*)]
  \item $T$ visits $A$ with frequency $p_A$ regardless of how small $p_A$
    is, and the visiting frequency entails no additional cost beyond the
    floor $\bmin$ already deposited by the partition.
  \item Locating $A$ (certifying that a given observation is in $A$ rather
    than in some $B \in \Sep(A,\Part)$) costs at least $\lambda(A) > 0$.
  \item As $p_A \to 0$, visiting frequency $\to 0$ but locating cost
    $\to +\infty$.  These are categorically distinct relations.
\end{enumerate}
\end{theorem}

\begin{proof}
(i) Follows directly from the Birkhoff ergodic theorem applied to
$\mathbf{1}_A$ \citep{halmos1950measure}.  The ergodic average converges to
$\mu(A) = p_A \cdot \mu(\Omega)$.  No additional information beyond
$\Part$ is required to perform the traversal; the partition was already
deposited at cost $\geq \bmin$.

(ii) To locate $A$, one must distinguish it from all $B \in \Sep(A,\Part)$.
Each binary decision requires at least $\bmin$ information (by
Theorem~\ref{thm:floor-positive}).  A binary search tree over the
$\lceil\log_2(1/p_A)\rceil$ bits needed to identify $A$ within the
partition structure therefore costs at least $\lambda(A) = \bmin \cdot
\lceil\log_2(1/p_A)\rceil$.

(iii) As $p_A \to 0$, the Birkhoff average converges to $0$ (the set is
visited rarely), while $\lceil\log_2(1/p_A)\rceil \to +\infty$ (exponentially
many partition decisions are needed to isolate it).  These are formally
distinct: no re-parameterisation of one yields the other.
\end{proof}

\subsection{Non-Self-Verifiability}

The incompletability of individuation has an immediate consequence: no finite
agent can fully verify its own individuation status.

\begin{proposition}[Non-Self-Verifiability]
\label{prop:non-self-verify}
Let $\alpha \in \Part$ be an agent performing partition operations on $\BRS$.
Then $\alpha$ cannot certify, by any finite sequence of partition operations,
that its own individuation is complete.
\end{proposition}

\begin{proof}
To certify that $\alpha$ is individuated requires completing the comparison
of $\alpha$ against $\Neg(\alpha, \Part)$.  Since $\alpha$ itself is within
$\Omega$, and any partition operation performed by $\alpha$ is itself an
individuation act (a cell carved from $\Neg(\alpha, \Part)$), the
verification sub-problem is of the same kind as the original: each
verification step requires a new partition operation, which in turn requires
verification.  By Theorem~\ref{thm:incompletable}, no finite sequence of
operations completes the comparison.  Hence $\alpha$ cannot certify its own
individuation by any finite procedure.
\end{proof}

\begin{remark}
The irreducible residue $\bmin$ is precisely the measure of the minimum cell
that any single individuation step must deposit.  It cannot be exhibited as
a locatable set within the verifier's own partition, because locating it
would require a further individuation of that cell --- which deposits another
cell of measure $\geq \bmin$, and so on.  This is the measure-theoretic
analogue of G\"{o}del's incompleteness \citep{godel1931formal}: the floor
is a true property of the space that cannot be exhibited within any finite
partition generated by an agent within that space.  The parallel is
structural, not metaphorical: both involve a system attempting to certify
a property of itself, encountering an irreducible residue at each stage.
\end{remark}

\subsection{Non-Return Theorem}

\begin{definition}[Committed Record]
\label{def:committed-record}
Let $(\Part^{(t)})_{t \geq 0}$ be a sequence of partitions produced by a
contact map investigation.  The \emph{committed record} at time $t$ is
\begin{equation}
  M(t) \;=\; \sum_{s=0}^{t} \sum_{A \in \Part^{(s)}} \mathbf{1}[\bP^{(s)}(A)
    \text{ is committed}],
\end{equation}
the total count of partition depths that have been committed (written to a
permanent record).
\end{definition}

\begin{theorem}[Non-Return]
\label{thm:non-return}
The committed record $M(t)$ is monotone non-decreasing: $M(t+1) \geq M(t)$
for all $t \geq 0$.  No investigation process can return to a prior
individuation state once a partition depth has been committed.
\end{theorem}

\begin{proof}
Commitment is defined as the operation of writing a partition depth to a
permanent record.  A permanent record, by definition, is a function
$R : \Nat \to \Real_{\geq 0}$ with $R(t+1) \geq R(t)$ for all $t$ (the
record accumulates, never deletes).  Hence $M(t) = |R^{-1}((0,\infty)) \cap
[0,t]|$ is non-decreasing.  The non-return property follows: no finite
process can reduce $M$, because doing so would require deleting a committed
record entry, which by definition is impossible.
\end{proof}

\begin{remark}
The non-return theorem is the measure-theoretic signature of the
Incompletability Theorem (Theorem~\ref{thm:incompletable}): because the
comparison against $\Neg(\alpha, \Part)$ never terminates, the committed
record can only grow.  Each committed individuation advances the record
irreversibly.  This gives the contact map investigation a strict arrow of
time given by the accumulation of irreducible residues --- analogous to the
second law of thermodynamics for investigation.
\end{remark}

% ==========================================================================
\section{Truth as a Cell: The Empirical Structure of Individuation}
\label{sec:empirical}
% ==========================================================================

The mathematical framework of the preceding section is not an abstraction
imposed on the world.  It is read off from the actual structure of scientific
practice.  This section establishes the empirical grounding of the framework
through a single, decisive observation: \emph{if truth were a point, one
instrument at sufficient resolution would terminate any investigation.
Investigations never terminate.}

\subsection{The Resolution Argument}

The history of measurement is a history of increasing resolution.
Lithium --- element 3 --- was first characterised by its spectrum (1817),
then by its atomic mass (Berzelius, 1818), then by its electron configuration
(quantum mechanics, 1920s), then by its nuclear spin states (NMR, 1950s),
then by its crystal structure at sub-angstrom resolution (synchrotron
diffraction, 1990s), then by its isotopic fractionation patterns (mass
spectrometry, 2000s), then by its quantum-mechanical tunnelling behaviour
in lithium-ion battery anodes (first-principles DFT, 2010s).

At no point did any measurement terminate the investigation.  At each stage,
a new instrument accessed a previously invisible dimension of the same object,
and the investigation continued.

If truth were a point --- a zero-measure entity in measurement space --- then
one instrument at sufficient resolution would place a pointer exactly on it,
and the investigation would be complete.  The pointer has never landed.  The
investigation has never completed.  This is not contingent on the state of
technology; it is structural.

\begin{proposition}[Truth is a Cell]
\label{prop:truth-cell}
The \emph{truth} about any individuated item $X$ is not a point in a
measurement space but a \emph{cell}: the minimum realisable unit of a
partition that carries positive measure and is named by the partition scheme.
Equivalently, the truth about $X$ is the partition floor $\bmin$ of the
individuation of $X$ against $\Neg(X, \Part)$.
\end{proposition}

\begin{proof}
A point has zero measure.  By Definition~\ref{def:detectability}, a region
of zero measure is subdetectable: it is engulfed by its own boundary and
cannot be distinguished from adjacent regions by any measure-preserving
observation.  If truth were a point, it would be subdetectable --- it would
be indistinguishable from the boundary between ``true'' and ``not true,''
which is not a truth.

A cell, by Definition~\ref{def:partition}(iv), has positive measure.  By
Theorem~\ref{thm:floor-positive}, the minimum cell in any valid partition
has measure $\geq \bmin > 0$.  A cell is detectable: it can be distinguished
from adjacent cells by a measure-preserving observation with positive
probability.  The minimum sufficient cell --- the contact ground state of
truth --- is the partition floor $\bmin$.
\end{proof}

\begin{remark}
This does not mean that truth is imprecise or subjective.  It means that
truth is a \emph{region} of positive measure, not a zero-measure point.
The region is sharp: it has a definite boundary, a definite measure, and a
definite relation to its complement.  But no instrument can compress it to
zero without ceasing to be an instrument (a zero-measure region is
undetectable and therefore not a measurement).
\end{remark}

\subsection{The Separator is Knowledge}

\begin{proposition}[Knowledge as Separator Measure]
\label{prop:knowledge-separator}
Let $X$ be an individuated item and $\Part$ a partition containing $X$.
Everything known about $X$ --- every certified fact, every measurement,
every observation --- is the committed portion of the separator
$\Sep(X, \Part)$.  The total measure $\mu(\Sep(X, \Part))$ is
simultaneously:
\begin{enumerate}[label=(\roman*)]
  \item the total committed knowledge about $X$ (the sum of all partition
    operations certified against $\Neg(X, \Part)$), and
  \item the lower bound on the remaining unknown (the partition floor
    $\bmin$, the minimum measure any further partition operation must deposit).
\end{enumerate}
These are not two different quantities.  They are the same quantity viewed
from two directions: the separator is the knowledge, and the separator's
incompleteness is the irreducible unknown.
\end{proposition}

\begin{proof}
(i) Each partition operation $\mathcal{O}^{(j)}$ carves a cell $C_j$ from
$\Neg(X, \Part)$ and deposits it into the separator (certifying ``$X$ is
not $C_j$'').  The total separator after $n$ operations is
$\bigcup_{j=1}^n C_j$.  The measure $\sum_j \mu(C_j)$ is the cumulative
information deposited.

(ii) By Theorem~\ref{thm:incompletable}, the comparison against
$\Neg(X, \Part)$ is never complete.  The residual $R = \Neg(X, \Part)
\setminus \bigcup_j C_j$ always has $\mu(R) > 0$.  Each future operation
must deposit at least $\bmin$ into $R$.  Hence $\bmin$ is simultaneously
the lower bound on the cost of any further knowledge and the measure of
the irreducible unknown.
\end{proof}

\subsection{The Lithium Example: An Instrument Array as a Contact Map}

Let $X = \text{Lithium}$.  The following table lists the major partition
operations performed on Lithium and the dimension of S-entropy space each
accesses.

\begin{center}
\begin{tabular}{lll}
  \toprule
  Instrument & Partition operation & S-entropy dimension \\
  \midrule
  Flame spectrometry & Emission spectrum & Photon energy \\
  Mass spectrometry  & Isotopic abundance & Mass/charge \\
  NMR spectroscopy   & Nuclear spin resonance & Spin-3/2 (${}^7$Li) \\
  X-ray diffraction  & Crystal periodicity & Spatial frequency \\
  First-principles DFT & Electronic structure & Wavefunction density \\
  Electrochemical cycling & Ion transport & Diffusion coefficient \\
  \bottomrule
\end{tabular}
\end{center}

Each instrument is a partition operation $\mathcal{O}_j : \Omega \to
\hat{\CM}$ accessing a different coordinate direction of the S-entropy
space $S(X) \in [0,1]^p$.  The contact map $\CMbmin(\text{Lithium},
\text{not-Lithium})$ is the synthesis of all these dimensions
simultaneously.  No single instrument produces it.  The array collectively
approximates it from below, monotonically, without ever reaching it.

\begin{proposition}[The Instrument Array is a Contact Map]
\label{prop:instruments-contact-map}
Let $\{\mathcal{O}_1, \ldots, \mathcal{O}_m\}$ be a collection of
measurement instruments applied to an individuated item $X$.  Their combined
output is a partial contact map estimate:
\begin{equation}
  \hat{\CM}_{\text{array}} \;=\;
    \bigwedge_{j=1}^m \mathcal{O}_j(X),
\end{equation}
where $\wedge$ is pointwise minimum (the most conservative estimate across
all instruments).  This estimate satisfies:
\begin{enumerate}[label=(\roman*)]
  \item $\hat{\CM}_{\text{array}} \geq \CMbmin(X, \Neg(X, \Part))$
    (from below: no instrument can go below the floor);
  \item Adding a new instrument $\mathcal{O}_{m+1}$ produces
    $\hat{\CM}' \leq \hat{\CM}_{\text{array}}$ (refinement: the estimate
    decreases toward $\CMbmin$, never increases);
  \item $\lim_{m \to \infty} \hat{\CM}_{\text{array}} = \CMbmin$
    (asymptotic completeness: the array converges to the floor as the
    number of independent instruments grows without bound).
\end{enumerate}
\end{proposition}

\subsection{Resolution Increasing is Region Reduction, Not Point Approach}

\begin{proposition}[Resolution Increase = Cell Subdivision]
\label{prop:resolution-cell}
When a finer instrument $\mathcal{O}'$ is applied to $X$ in place of a
coarser instrument $\mathcal{O}$, the effect is a cell subdivision: a cell
$C \in \Part$ of measure $m$ is replaced by $C_1, C_2$ with $\mu(C_1) +
\mu(C_2) = m$ and $\mu(C_1), \mu(C_2) > 0$.  The floor descends:
$\bmin' \leq \bmin$.  But $\bmin' > 0$ (both subcells have positive
measure, by Definition~\ref{def:partition}(iv)).

A sequence of resolutions $\delta_1 > \delta_2 > \cdots \to 0$ is the
refinement chain of the BRS axioms (condition~\eqref{eq:resolution-floor});
the limit $\delta = 0$ is excluded for any finite agent.  Resolution
increasing without bound is the observable signature of individuation by
negation against an infinite complement.  It is \emph{reduction of a
region}, not \emph{approach to a point}.
\end{proposition}

\begin{remark}
This proposition makes the abstract framework concrete.  A mass spectrometer
that distinguishes ${}^6\text{Li}$ from ${}^7\text{Li}$ (resolving power
$m/\Delta m \approx 1000$) performs a cell subdivision of the ``Lithium
isotope'' cell.  A high-resolution instrument ($m/\Delta m \approx 100{,}000$)
subdivides further.  Neither reaches a point.  Each produces two subcells of
positive measure.  The floor descends with each subdivision but never
reaches zero.  The ``truth'' about Lithium's isotopic composition is the
current finest cell, which is always a region, always positive in measure,
always further subdivisible in principle.
\end{remark}

\subsection{Invariance of the Floor}

\begin{theorem}[Water Level Invariance]
\label{thm:water-beads}
Let $\Sigma = \{\sigma_1, \ldots, \sigma_N\}$ be a finite set of
\emph{observable states} (test outcomes, type annotations, measurement
results, runtime events).  An \emph{arrangement} $\alpha : \Part \to
2^\Sigma$ assigns states to components.  For any rearrangement
$\alpha' = \pi \circ \alpha$ (where $\pi$ is any permutation of state
assignments):
\begin{equation}
  \bmin(\alpha) \;=\; \bmin(\alpha') \;=\; \bmin.
\end{equation}
The partition floor is invariant under all rearrangements of observable
states.
\end{theorem}

\begin{proof}
The partition floor $\bmin$ is a property of the measure space $\BRS$ and
the partition $\Part$; it does not depend on the assignment $\alpha$.
Formally: $\bmin = \inf_\Part \inf_{A \in \Part} \inf_{B \in \Sep(A,\Part)}
\mu(B)$.  The measure $\mu$ is a fixed Borel measure on $\Omega$; neither
$\alpha$ nor $\pi$ acts on $\mu$ or on the topological structure of $\Omega$.
Hence $\bmin(\alpha) = \bmin(\alpha') = \bmin$.
\end{proof}

\begin{remark}
This formalises the \emph{beads and water} intuition.  Observable states
$\Sigma$ (test results, instrument readings, type errors) are the beads:
their arrangement and rearrangement does not change the water level $\bmin$.
A system with $N$ passing tests and a system with $N$ failing tests have
the same partition floor if they have the same underlying bounded resolvable
space.  The water level is the sufficient amount of unknowns: it is
constituted by the incompletable comparison against the complement, not by
the arrangement of what has been found.
\end{remark}

\subsection{Software Implication}

The software consequence follows immediately.  A type checker is one
instrument.  A unit test suite is another.  A static analyser is another.
A runtime trace is another.  An LLM inferring semantic boundaries is another.
Each is a partition operation in a different coordinate direction of the
S-entropy space of the software system.  None alone produces $\CMbmin$.

The contact map generator is the instrument array --- synthesising all
partition operations into a running estimate of the contact ground state.
The \texttt{compose} block in the Wind Tunnel DSL is the synthesis step
(pointwise minimum across instruments).  The \texttt{recurse} block is the
refinement chain (subdividing cells when marginal return justifies it).  The
termination condition \texttt{until marginal\_return < floor} is literally
$\bmin$: the rational halt of a finite agent in an infinite incompletable
comparison.  The halt is not failure --- it is the correct recognition that
the floor has been reached and further operations would deposit less than one
cell's worth of new information.

% ==========================================================================
\section{The Minimum Sufficient Contact Map}
\label{sec:mscm}
% ==========================================================================

\subsection{S-Entropy Coordinates}

\begin{definition}[S-Entropy of a Component]
\label{def:s-entropy}
Let $\Part \in \mathfrak{P}(\Omega,d,\mu)$ and $A \in \Part$.  Fix a
$p$-dimensional feature map $\phi : \Omega \to [0,1]^p$.  The
\emph{S-entropy coordinate vector} of $A$ is
\begin{equation}
  \Scoord(A) \;=\; \left( H_\phi^{(1)}(A), \ldots, H_\phi^{(p)}(A) \right)
    \;\in\; [0,1]^p,
\end{equation}
where
\begin{equation}
  H_\phi^{(j)}(A) \;=\; -\int_A \phi_j(x) \log_2 \phi_j(x)\, d\mu(x)
    \;\bigg/\; \mu(A)
\end{equation}
is the normalised Shannon entropy of the $j$-th feature within $A$
\citep{shannon1948mathematical, cover2006elements}.  The \emph{S-entropy
distance} between two components $A_i, A_j \in \Part$ is
\begin{equation}
  \dS(A_i, A_j) \;=\; \left\|\Scoord(A_i) - \Scoord(A_j)\right\|_2.
\end{equation}
\end{definition}

\begin{remark}
The S-entropy coordinates capture the distributional geometry of each
component in feature space.  $\dS$ is a \emph{pseudo-metric} on $\Part$:
it satisfies non-negativity, symmetry, and the triangle inequality, but
$\dS(A_i, A_j) = 0$ does not imply $A_i = A_j$ in general.  It is a
genuine metric when the feature map $\phi$ is \emph{sufficiently
discriminative}: $\phi(A_i) \neq \phi(A_j)$ whenever $A_i$ and $A_j$ are
distinct components.  All theorems below that require strict separation
between components implicitly require this discriminativity condition.
\end{remark}

\subsection{The Contact Graph and Contact Map}

\begin{definition}[Contact Graph]
\label{def:contact-graph}
The \emph{contact graph} of $\Part$ is the undirected graph
$\GC = (\Part, E_C)$ where
\begin{equation}
  E_C = \left\{\{A_i, A_j\} : A_j \in \Sep(A_i, \Part)\right\}.
\end{equation}
Two components are adjacent in $\GC$ if and only if their closures share a
boundary point.
\end{definition}

\begin{definition}[Contact Map]
\label{def:contact-map}
The \emph{contact map} of $(\Part, \phi)$ is the edge-weighting
\begin{equation}
  \CM : E_C \;\longrightarrow\; \Real_{\geq 0}, \qquad
  \CM(A_i, A_j) \;=\; \dS(S(A_i), S(A_j)).
\end{equation}
(When $\phi$ is not fully discriminative, $\CM$ is a pseudo-metric on $E_C$.)
\end{definition}

\subsection{The Minimum Sufficient Contact Map}

\begin{definition}[Contact Locus]
\label{def:contact-locus}
The \emph{contact locus} of $\Part$ is the sub-collection
\begin{equation}
  \Part_{\bmin} \;=\; \left\{ A \in \Part : \bP(A) = \bmin \right\}
\end{equation}
of components that sit precisely at the partition floor.  The
\emph{minimum sufficient contact map} is the restriction
\begin{equation}
  \CMbmin(A_i, A_j) \;=\; \dS\!\left(\Sbmin(A_i),\, \Sbmin(A_j)\right)
\end{equation}
for $(A_i, A_j) \in E_C$ with $A_i, A_j \in \Part_{\bmin}$, where
$\Sbmin(A) = \Scoord(A)|_{\Part_{\bmin}}$ are the S-entropy coordinates
evaluated at the contact locus.
\end{definition}

\begin{definition}[Partition Depth Inflation]
\label{def:inflation}
A contact map $\CM'$ is a \emph{partition depth inflation} of $\CMbmin$ if
there exists a non-negative functional $\delta\CM \geq 0$ such that
\begin{equation}
  \CM'(A_i, A_j) \;=\; \CMbmin(A_i, A_j) + \delta\CM(A_i, A_j)
\end{equation}
for all $(A_i, A_j) \in E_C$.  Every inflated contact map is further from
the contact ground state than $\CMbmin$.
\end{definition}

\subsection{The Minimum Sufficient Lagrangian}

\begin{definition}[Partition Lagrangian]
\label{def:lagrangian}
Given a path $\gamma : [0,T] \to \Omega$ in the bounded resolvable space,
the \emph{partition Lagrangian} associated to a contact map $\CM$ is
\begin{equation}
  L(\gamma, \dot\gamma, t) \;=\; \tfrac{1}{2}\,\mu\!\left(B(\gamma(t), \|\dot\gamma(t)\|)\right)
    \;-\; M_\CM(\gamma(t)),
\end{equation}
where $M_\CM : \Omega \to \Real$ is the \emph{contact potential} induced by
$\CM$:
\begin{equation}
  M_\CM(x) \;=\; \sum_{(A_i,A_j) \in E_C} \CM(A_i,A_j)\,
    \mathbf{1}_{A_i}(x)\,\mathbf{1}_{A_j}(x).
\end{equation}
The \emph{minimum sufficient Lagrangian} is $\Lmin = L|_{\CM = \CMbmin}$,
corresponding to the contact potential $M_{\bmin} = M_\CM|_{\CM = \CMbmin}$.
\end{definition}

\subsection{The Minimum Sufficiency Trijection}

\begin{theorem}[Minimum Sufficiency Trijection]
\label{thm:minimum-sufficiency}
The following three objects are in canonical bijective correspondence:
\begin{enumerate}[label=(\Roman*)]
  \item The partition floor $\bmin$;
  \item The minimum sufficient contact map $\CMbmin$;
  \item The minimum sufficient Lagrangian $\Lmin$.
\end{enumerate}
Every other configuration of $\BRS$ is a partition depth inflation of
$\CMbmin$: the general contact map satisfies $\CM = \CMbmin + \delta\CM$
with $\delta\CM \geq 0$ pointwise, and the corresponding Lagrangian satisfies
$L = \Lmin + \delta L$ with $\delta L = -M_{\delta\CM} \leq 0$ pointwise
(inflation increases the contact potential, reducing the kinetic dominance
of $L$).
\end{theorem}

\begin{proof}
We establish the bijections by constructing explicit maps in each direction.

\textbf{(I) $\Rightarrow$ (II).}  Given $\bmin$, the contact locus
$\Part_{\bmin}$ is determined by Definition~\ref{def:contact-locus}.  The
S-entropy coordinates $\Sbmin(A)$ for $A \in \Part_{\bmin}$ are determined
by $\mu$ and $\phi$ (fixed by the space).  Hence $\CMbmin$ is uniquely
determined by $\bmin$ (given a fixed, sufficiently discriminative $\phi$).

\textbf{(II) $\Rightarrow$ (I).}  Given $\CMbmin$, the contact locus
$\Part_{\bmin}$ is its domain.  By Definition~\ref{def:partition-depth},
$\bmin = \inf_{A \in \Part_{\bmin}} \bP(A)$.  Since $\CMbmin$ is defined
precisely on the contact locus, $\bmin$ is determined.

\textbf{(II) $\Rightarrow$ (III).}  Given $\CMbmin$, the contact potential
$M_{\bmin}$ is determined by Definition~\ref{def:lagrangian}, and hence
$\Lmin$ is uniquely determined.

\textbf{(III) $\Rightarrow$ (II).}  Given $\Lmin$, the contact potential
$M_{\bmin}$ is the negation of the non-kinetic term of $\Lmin$.  The
contact map values $\CMbmin(A_i,A_j)$ are the coefficients in $M_{\bmin}$,
hence determined.

For the inflation direction: a partition depth inflation adds $\delta\CM
\geq 0$ to the contact map, which adds $M_{\delta\CM} \geq 0$ to the
contact potential.  The Lagrangian then becomes
$L = \tfrac{1}{2}\mu(B(\gamma,\|\dot\gamma\|)) - (M_{\bmin} + M_{\delta\CM})
= \Lmin - M_{\delta\CM} \leq \Lmin$ pointwise.  Hence inflation reduces the
Lagrangian pointwise.  The contact map distance $\CM(A_i,A_j) =
\CMbmin(A_i,A_j) + \delta\CM(A_i,A_j) \geq \CMbmin(A_i,A_j)$ --- inflation
moves the contact map \emph{away} from the ground state, not toward it.
\end{proof}

\begin{corollary}[Minimum Sufficient Action]
\label{cor:action}
The \emph{minimum sufficient action} is
\begin{equation}
  \Smin \;=\; \int_0^T \Lmin(\gamma, \dot\gamma, t)\, dt.
\end{equation}
It is determined by $\Lmin$ via time integration and therefore by
$\CMbmin$ via the trijection.  It is \emph{not} a fourth independent
object: distinct Lagrangians can yield the same action (e.g., differing by
boundary terms or total-derivative additions), so $\Smin$ does not determine
$\Lmin$ uniquely.  The minimum sufficient action is a derived quantity, not
a member of the trijection.
\end{corollary}

\begin{corollary}[Contact Ground State]
\label{cor:ground-state}
The minimum sufficient contact map $\CMbmin$ is the \emph{contact ground
state}: the minimum structure required to establish that every component is
surely not any other component, against the whole system.  Every additional
structural annotation (type declarations, test suites, runtime traces,
lint rules) is a partition depth inflation $\delta\CM \geq 0$ of this
ground state.  None can reduce the floor to zero; all are upward
perturbations of the contact map distance.
\end{corollary}

% ==========================================================================
\section{Software Systems as Bounded Resolvable Spaces}
\label{sec:software}
% ==========================================================================

\subsection{The System Graph}

\begin{definition}[Software System]
\label{def:software-system}
A \emph{software system} is a quadruple $\mathfrak{S} = (U, E, A, \Pi)$
where:
\begin{enumerate}[label=(\roman*)]
  \item $U = \{u_1, \ldots, u_N\}$ is a finite set of \emph{units} (functions,
    classes, modules, services, or any other named computation boundary);
  \item $E \subseteq U \times U$ is a set of directed \emph{edges}
    representing invocation, data-flow, or dependency relations;
  \item $A : U \to 2^\mathcal{T}$ assigns to each unit an \emph{aperture set}
    --- the collection of types it consumes;
  \item $\Pi : U \to 2^\mathcal{T}$ assigns to each unit a
    \emph{production set} --- the collection of types it produces;
\end{enumerate}
where $\mathcal{T}$ is a fixed type universe.  The \emph{system graph}
$\GS = (U, E)$ is the underlying directed graph.
\end{definition}

\begin{definition}[Software Metric Space]
\label{def:software-metric}
Given $\mathfrak{S}$, define a metric on $U$ by
\begin{equation}
  d_\mathfrak{S}(u_i, u_j) \;=\;
    1 - \frac{|A(u_i) \cap A(u_j)| + |\Pi(u_i) \cap \Pi(u_j)|}
             {|A(u_i) \cup A(u_j)| + |\Pi(u_i) \cup \Pi(u_j)|},
\end{equation}
the complement of the Jaccard similarity of the combined aperture/production
type sets.  Define a measure on $U$ by $\mu(u_i) = w_i$ where $w_i$ is the
normalised coupling weight (e.g., proportional to in-degree in $\GS$).
\end{definition}

\begin{proposition}
$(U, d_\mathfrak{S}, \mu)$ is a bounded resolvable space (in the finite,
discrete sense) with resolution data $(1, w_{\min}, 1)$, where $w_{\min} =
\min_i w_i > 0$.
\end{proposition}

\begin{proof}
$U$ is finite, hence compact.  $\mu$ is a finite measure with $\mu(u_i) =
w_i > 0$ for all $i$ (by assumption of positive coupling weight).  The
resolution floor holds trivially: $\mu(B(u_i, r)) \geq w_i \geq w_{\min}$
for $r \geq 0$.  In this discrete finite setting, $\bmin = w_{\min} > 0$
follows directly: the minimum cell measure in any finite partition with
positive-weight cells is a minimum over a finite set of positive numbers,
hence positive.  This is the direct finite-graph version of
Theorem~\ref{thm:floor-positive}.
\end{proof}

\subsection{Holonomy of System Cycles}

\begin{definition}[Cycle Holonomy]
\label{def:holonomy}
Let $c = (u_{i_1}, u_{i_2}, \ldots, u_{i_k}, u_{i_1})$ be a directed cycle
in $\GS$.  The \emph{specification type flow} along $c$ is the sequence of
type transitions
\begin{equation}
  \tau(c) \;=\; \Pi(u_{i_1}) \to A(u_{i_2}) \to \Pi(u_{i_2}) \to \cdots
    \to A(u_{i_1}).
\end{equation}
The \emph{holonomy} of $c$ is
\begin{equation}
  \hol(c) \;=\; \left\| \tau(c) - \tau^{\mathrm{spec}}(c) \right\|,
\end{equation}
where $\tau^{\mathrm{spec}}(c)$ is the specified type flow and $\|\cdot\|$
is a norm on the space of type-transition sequences.
\end{definition}

\begin{remark}
Holonomy measures how much the actual type flow around a cycle deviates from
its specification.  A system with $\hol(c) = 0$ for all cycles is
\emph{specificationally consistent} around all loops.  Non-zero holonomy
indicates a coordination failure: the types produced by traversing the cycle
do not match the types expected.
\end{remark}

\subsection{S-Entropy of Software Units}

\begin{definition}[S-Entropy of a Software Unit]
\label{def:s-entropy-software}
For a software unit $u \in U$, fix a $p$-dimensional feature map
$\phi_\mathfrak{S} : U \to [0,1]^p$ encoding structural properties
(e.g.\ cyclomatic complexity, coupling degree, type-set cardinality,
dependency depth).  The \emph{S-entropy coordinate vector} of $u$ is
\begin{equation}
  \Scoord(u) \;=\; \left(-\phi_j(u)\log_2\phi_j(u)\right)_{j=1}^p
    \;\in\; [0,1]^p.
\end{equation}
The \emph{S-entropy} of $u$ is $S(u) = \|\Scoord(u)\|_1 \in [0, p\log_2 e]$.
The floor $S_{\flat} = \min_{u \in U} S(u)$ is strictly positive for any
non-trivially structured system.
\end{definition}

\subsection{Five Coordination Regimes}

\begin{definition}[Static Order Parameter]
\label{def:order-parameter}
The \emph{static order parameter} of $\mathfrak{S}$ is
\begin{equation}
  \Rest \;=\; \exp\!\left(-\,\overline{\tau}\right),
\end{equation}
where $\overline{\tau}$ is the mean tension:
\begin{equation}
  \overline{\tau} \;=\; \frac{1}{|E_C|} \sum_{(u_i, u_j) \in E_C}
    \dS(S(u_i), S(u_j)) \cdot \frac{\hol(c_{ij})}{\hol_{\max} + \varepsilon},
\end{equation}
with $c_{ij}$ the shortest cycle containing the edge $(u_i,u_j)$,
$\hol_{\max} = \max_c \hol(c)$, and $\varepsilon > 0$ a regulariser.
\end{definition}

\begin{definition}[Coordination Regimes]
\label{def:coordination-regimes}
The coordination regimes of a software system are determined by $\Rest$:

\begin{center}
\begin{tabular}{lll}
  \toprule
  Regime & $\Rest$ range & Characterisation \\
  \midrule
  Turbulent & $[0, 0.3)$
    & Maximum partition inflation; near-zero contact coherence \\
  Aperture-dominated & $[0.3, 0.5)$
    & Type boundaries dominate; structural coupling weak \\
  Hierarchical cascade & $[0.5, 0.8)$
    & Layered dependency structure; partial contact coherence \\
  Coherent & $[0.8, 0.95)$
    & High contact coherence; near minimum sufficient structure \\
  Phase-locked & $[0.95, 1]$
    & Near contact ground state; $\CM \approx \CMbmin$ \\
  \bottomrule
\end{tabular}
\end{center}

\textbf{Note on thresholds.}  The boundaries $0.3, 0.5, 0.8, 0.95$ are
\emph{operational parameters}, empirically calibrated against known
coordination failure modes in software systems.  They are not consequences
of the mathematical theory; they are quantitative demarcations within the
continuous range $[0,1]$ of $\Rest$ that have been found to correspond to
qualitative changes in system behaviour.  The regime classification should
be understood as a practical instrument, not a derived partition.
\end{definition}

\begin{remark}
$\Rest$ is an order parameter in the statistical mechanics sense: it measures
proximity of the system's actual contact map to the minimum sufficient contact
map.  $\Rest \to 1$ corresponds to $\CM \to \CMbmin$; $\Rest \to 0$
corresponds to maximally inflated partition depth ($\CM \gg \CMbmin$).
\end{remark}

\subsection{Local Invisibility Theorem}

\begin{theorem}[Local Invisibility]
\label{thm:local-invisibility}
Let $\mathfrak{S} = (U, E, A, \Pi)$ and let $\val_u : \mathcal{P}_u \to
\{0,1\}$ be a local verification predicate for unit $u$, where $\mathcal{P}_u$
is the set of properties of $u$ in isolation.  Then:
\begin{equation}
  \forall u \in U : \val_u(s) = 1 \;\not\Rightarrow\;
    \CM_\mathfrak{S} = \CMbmin^\mathfrak{S}.
\end{equation}
That is, universal local validity does not imply that the system's contact
map is at the minimum sufficient ground state.
\end{theorem}

\begin{proof}
Construct a counterexample.  Let $U = \{u_1, u_2, u_3\}$ with edges
$E = \{(u_1,u_2),(u_2,u_3),(u_3,u_1)\}$ (a 3-cycle).  Define aperture and
production sets so that each $u_i$ satisfies its local type contract in
isolation: $\Pi(u_i) \subseteq A(u_{i+1 \mod 3})$.  Now insert a coercion
$\kappa : T_1 \to T_1'$ that is invisible to each local verifier (e.g., a
silent widening conversion in a language with subtype coercion).  Each
$\val_{u_i}(s) = 1$ because the local type contract is satisfied with the
coercion.  However, the holonomy of the cycle is $\hol(c) = \|\kappa\| > 0$:
after one traversal of the cycle, the type has widened silently, deviating
from the specification.  Hence $\CM_\mathfrak{S}(u_i, u_j) >
\CMbmin^\mathfrak{S}(u_i, u_j)$ for some $(u_i,u_j)$ --- the contact map
is inflated above the ground state despite universal local validity.

The general impossibility follows from Rice's theorem \citep{rice1953classes}:
the global semantic property $\CM_\mathfrak{S} = \CMbmin$ is non-trivial and
hence undecidable from local information alone.  Local predicates $\val_u$
capture only decidable local properties and are therefore insufficient to
determine it.
\end{proof}

% ==========================================================================
\section{The Contact Map Generator}
\label{sec:generator}
% ==========================================================================

\subsection{What the Contact Map Generator Is}

The \emph{Contact Map Generator} (CMG) is a system that, given a software
system $\mathfrak{S}$, produces a sequence of progressively refined estimates
\begin{equation}
  \hat{\CM}^{(0)} \;\leq\; \hat{\CM}^{(1)} \;\leq\; \cdots \;\leq\; \hat{\CM}^{(T)}
    \;\xrightarrow{T\to\infty}\; \CMbmin^\mathfrak{S}
\end{equation}
of the minimum sufficient contact map, where $\leq$ denotes refinement in the
sense of decreasing partition depth inflation ($\delta\CM$ monotone
non-increasing).

The CMG is not a test runner.  It is not a linter.  It is not a type checker.
It is a \emph{partition machine}: it accepts any process --- static analysis,
dynamic tracing, LLM-based semantic inference, user-provided structural
intuition --- as a \emph{partition operation} that deposits information about
$\CMbmin^\mathfrak{S}$.  Each operation is one instrument in the array;
the CMG is the array itself.

\subsection{Partition Operations}

\begin{definition}[Partition Operation]
\label{def:partition-operation}
A \emph{partition operation} on $\mathfrak{S}$ is a computable function
\begin{equation}
  \mathcal{O} : \mathfrak{S} \;\longrightarrow\; \hat{\CM}
\end{equation}
producing a partial contact map estimate $\hat{\CM}$ (a partial function on
$E_C$).  Partition operations include:
\begin{enumerate}[label=(\alph*)]
  \item \textbf{Static}: call graph extraction, type aperture analysis,
    dependency edge analysis \citep{nielson1999principles, aho2006compilers};
  \item \textbf{Dynamic}: runtime trace analysis, coverage profiling,
    execution-path sampling;
  \item \textbf{Semantic}: LLM-based inference of coordination boundaries
    from code text;
  \item \textbf{Structural}: user-provided assertions about architectural
    boundaries.
\end{enumerate}
\end{definition}

\begin{definition}[Composition of Partition Operations]
\label{def:composition}
Given two partition operations $\mathcal{O}_1, \mathcal{O}_2$ producing
estimates $\hat{\CM}_1, \hat{\CM}_2$, their \emph{composition} is
\begin{equation}
  (\mathcal{O}_1 \oplus \mathcal{O}_2)(\mathfrak{S}) \;=\;
    \hat{\CM}_1 \wedge \hat{\CM}_2,
\end{equation}
where $\wedge$ takes the pointwise minimum (the more conservative estimate
on each edge), respecting the monotone non-inflation property.
\end{definition}

\subsection{The Cessation Theorem}

\begin{theorem}[Cessation]
\label{thm:cessation}
Let $(\mathcal{O}^{(t)})_{t \geq 0}$ be a sequence of partition operations
applied by the CMG, with marginal returns $\Delta_t = \|\hat{\CM}^{(t)} -
\hat{\CM}^{(t-1)}\|$.  Then $\Delta_t \to \bmin > 0$ as $t \to \infty$;
the CMG halts (ceases further operations) when $\Delta_t < \bmin$.
\end{theorem}

\begin{proof}
By Theorem~\ref{thm:floor-positive}, $\bmin > 0$.  Each partition operation
$\mathcal{O}^{(t)}$ deposits at least $\bmin$ information (by
Definition~\ref{def:partition-depth} and the floor positivity), so $\Delta_t
\geq \bmin$ for all $t$ until the estimate converges to $\CMbmin$.  Once
$\hat{\CM}^{(t)} = \CMbmin$ (up to the resolution of the feature map $\phi$),
further operations produce $\Delta_t < \bmin$ (no new information above
the floor), and the cessation condition is met.  Hence the CMG halts in
finite time --- at or near the contact ground state.
\end{proof}

% ==========================================================================
\section{The Wind Tunnel DSL}
\label{sec:dsl}
% ==========================================================================

\subsection{Overview and Design Principles}

The \emph{Wind Tunnel DSL} is the language in which contact map
investigations are specified and recorded.  It is not a test language and not
a configuration language.  Its design principles are:

\begin{enumerate}[label=(\roman*)]
  \item \textbf{Partition-first}: every construct corresponds to a partition
    operation or a claim about partition depths.
  \item \textbf{Co-authored}: both users and the CMG system write valid DSL.
    The same syntax is used by both.
  \item \textbf{Self-describing}: a completed DSL execution file is a
    complete, reproducible record of a contact map investigation.
  \item \textbf{Monotone}: committed partition claims cannot be undone within
    a script (Non-Return Theorem, Theorem~\ref{thm:non-return}).
  \item \textbf{Floor-aware}: every \texttt{assert} is a minimum individuation
    claim (a lower bound on $\CMbmin$), never a truth-at-a-point claim.
\end{enumerate}

\subsection{Formal Grammar}

The Wind Tunnel DSL is defined by the following context-free grammar in BNF
form.

\begin{definition}[Wind Tunnel Grammar]
\label{def:grammar}
\begin{align*}
\langle\textit{program}\rangle &\;::=\; \langle\textit{block}\rangle^+ \\[4pt]
\langle\textit{block}\rangle &\;::=\;
  \langle\textit{scope-block}\rangle \;\mid\; \langle\textit{analyse-block}\rangle \\
  &\quad\;\mid\; \langle\textit{compose-block}\rangle \;\mid\;
    \langle\textit{assert-block}\rangle \\
  &\quad\;\mid\; \langle\textit{recurse-block}\rangle \;\mid\;
    \langle\textit{report-block}\rangle \\[4pt]
\langle\textit{scope-block}\rangle &\;::=\;
  \texttt{scope}\;\langle\textit{id}\rangle\;\texttt{\{}\;
  \langle\textit{scope-body}\rangle\;\texttt{\}} \\
\langle\textit{scope-body}\rangle &\;::=\;
  \langle\textit{target-decl}\rangle^+\;\langle\textit{filter-decl}\rangle^* \\
\langle\textit{target-decl}\rangle &\;::=\;
  \texttt{target:}\;\langle\textit{glob-pattern}\rangle \\
\langle\textit{filter-decl}\rangle &\;::=\;
  \texttt{exclude:}\;\langle\textit{glob-pattern}\rangle
  \;\mid\; \texttt{depth:}\;\langle\textit{nat}\rangle \\[4pt]
\langle\textit{analyse-block}\rangle &\;::=\;
  \texttt{analyse}\;\langle\textit{method}\rangle\;\texttt{\{}\;
  \langle\textit{analyse-body}\rangle\;\texttt{\}} \\
\langle\textit{method}\rangle &\;::=\;
  \texttt{static} \;\mid\; \texttt{dynamic} \;\mid\;
  \texttt{semantic} \;\mid\; \texttt{hf} \;\mid\; \langle\textit{id}\rangle \\
\langle\textit{analyse-body}\rangle &\;::=\;
  \langle\textit{step}\rangle^+ \\
\langle\textit{step}\rangle &\;::=\;
  \texttt{extract:}\;\langle\textit{id}\rangle
  \;\mid\; \texttt{model:}\;\langle\textit{string}\rangle
  \;\mid\; \texttt{query:}\;\langle\textit{string}\rangle
  \;\mid\; \texttt{method:}\;\langle\textit{id}\rangle \\[4pt]
\langle\textit{compose-block}\rangle &\;::=\;
  \texttt{compose}\;\texttt{\{}\;
  \langle\textit{compose-body}\rangle\;\texttt{\}} \\
\langle\textit{compose-body}\rangle &\;::=\;
  \langle\textit{merge-decl}\rangle^+ \\
\langle\textit{merge-decl}\rangle &\;::=\;
  \texttt{merge:}\;\texttt{[}\;\langle\textit{id-list}\rangle\;\texttt{]}
  \;\mid\; \texttt{resolve:}\;\langle\textit{resolve-strategy}\rangle \\
\langle\textit{resolve-strategy}\rangle &\;::=\;
  \texttt{overlaps by minimum\_depth}
  \;\mid\; \texttt{overlaps by maximum\_depth}
  \;\mid\; \langle\textit{id}\rangle \\[4pt]
\langle\textit{assert-block}\rangle &\;::=\;
  \texttt{assert}\;\langle\textit{assertion-expr}\rangle\;\texttt{\{}\;
  \langle\textit{assert-body}\rangle\;\texttt{\}} \\
\langle\textit{assertion-expr}\rangle &\;::=\;
  \langle\textit{individuation-claim}\rangle
  \;\mid\; \langle\textit{distance-claim}\rangle
  \;\mid\; \langle\textit{regime-claim}\rangle \\
\langle\textit{individuation-claim}\rangle &\;::=\;
  \texttt{individuation} \\
\langle\textit{distance-claim}\rangle &\;::=\;
  \texttt{distance} \\
\langle\textit{regime-claim}\rangle &\;::=\;
  \texttt{regime} \\
\langle\textit{assert-body}\rangle &\;::=\;
  \texttt{claim:}\;\langle\textit{claim-expr}\rangle
  \;\langle\textit{on-fail-decl}\rangle^? \\
\langle\textit{claim-expr}\rangle &\;::=\;
  \texttt{distance(}\langle\textit{id}\rangle\texttt{,}\langle\textit{id}\rangle\texttt{)}\;
  \texttt{>=}\;\langle\textit{float}\rangle \\
  &\quad\;\mid\; \texttt{R\_est}\;\texttt{>=}\;\langle\textit{float}\rangle
  \;\mid\; \langle\textit{id}\rangle\;\langle\textit{op}\rangle\;\langle\textit{float}\rangle \\
\langle\textit{on-fail-decl}\rangle &\;::=\;
  \texttt{on\_fail:}\;\langle\textit{fail-action}\rangle \\
\langle\textit{fail-action}\rangle &\;::=\;
  \texttt{narrow\_scope(}\langle\textit{id}\rangle\texttt{,}\langle\textit{id}\rangle\texttt{)}
  \;\mid\; \langle\textit{recurse-block}\rangle
  \;\mid\; \langle\textit{id}\rangle \\[4pt]
\langle\textit{recurse-block}\rangle &\;::=\;
  \texttt{recurse}\;\langle\textit{condition}\rangle\;\texttt{\{}\;
  \langle\textit{recurse-body}\rangle\;\texttt{\}} \\
\langle\textit{condition}\rangle &\;::=\;
  \texttt{until}\;\langle\textit{threshold-expr}\rangle
  \;\mid\; \texttt{while}\;\langle\textit{cond-expr}\rangle \\
\langle\textit{threshold-expr}\rangle &\;::=\;
  \texttt{marginal\_return <}\;\langle\textit{id}\rangle
  \;\mid\; \langle\textit{id}\rangle\;\langle\textit{op}\rangle\;\langle\textit{float}\rangle \\
\langle\textit{recurse-body}\rangle &\;::=\;
  \langle\textit{block}\rangle^+ \\[4pt]
\langle\textit{report-block}\rangle &\;::=\;
  \texttt{report}\;\texttt{\{}\;\langle\textit{report-body}\rangle\;\texttt{\}} \\
\langle\textit{report-body}\rangle &\;::=\;
  \langle\textit{output-decl}\rangle^+ \\
\langle\textit{output-decl}\rangle &\;::=\;
  \texttt{output:}\;\langle\textit{id}\rangle
  \;\mid\; \texttt{format:}\;\langle\textit{id}\rangle
\end{align*}
\end{definition}

\subsection{Operational Semantics}

We define a small-step operational semantics for the DSL.

\begin{definition}[DSL State]
\label{def:dsl-state}
A \emph{DSL state} is a triple $\Gamma = (\hat{\CM}, \mathcal{C}, \mathcal{R})$
where:
\begin{enumerate}[label=(\roman*)]
  \item $\hat{\CM}$ is the current contact map estimate;
  \item $\mathcal{C} \subseteq E_C$ is the set of committed edges (pairs
    whose partition depth has been certified);
  \item $\mathcal{R} \in \{\textit{pass}, \textit{fail}, \textit{narrow},
    \textit{recurse}\}$ is the current resolution flag.
\end{enumerate}
The initial state is $\Gamma_0 = (\hat{\CM}^{(0)}, \emptyset, \textit{pass})$
where $\hat{\CM}^{(0)}$ is the trivial (all-unknown) estimate.
\end{definition}

\begin{definition}[Block Reduction Rules]
\label{def:reduction-rules}
The reduction relation $\Gamma \xrightarrow{b} \Gamma'$ for each block
type $b$ is defined as follows.

\paragraph{Scope.}
$(\hat{\CM}, \mathcal{C}, \mathcal{R})
  \xrightarrow{\texttt{scope}(id, \textit{targets})}
  (\hat{\CM}, \mathcal{C}, \mathcal{R})[\textit{active-scope} := \textit{targets}]$.
The scope block declares the active partition region.  It does not modify the
contact map or commit any partition.

\paragraph{Analyse.}
$(\hat{\CM}, \mathcal{C}, \mathcal{R})
  \xrightarrow{\texttt{analyse}(\mathit{method}, \textit{steps})}
  (\hat{\CM} \wedge \hat{\CM}_{\mathit{method}}, \mathcal{C}, \mathcal{R})$,
where $\hat{\CM}_{\mathit{method}}$ is the partial contact map produced by
partition operation $\mathcal{O}_{\mathit{method}}$ applied to the active
scope.  The composition $\wedge$ is pointwise minimum.

\paragraph{Compose.}
$(\hat{\CM}, \mathcal{C}, \mathcal{R})
  \xrightarrow{\texttt{compose}(\textit{merge-list})}
  (\bigwedge_j \hat{\CM}_j, \mathcal{C}, \mathcal{R})$,
where $\hat{\CM}_j$ are the estimates named in the merge list and $\bigwedge$
is pointwise minimum over all named estimates.

\paragraph{Assert (success).}
If $\textit{claim}$ evaluates to $\top$ in state $\Gamma$:
$(\hat{\CM}, \mathcal{C}, \mathcal{R})
  \xrightarrow{\texttt{assert}(\textit{claim})}
  (\hat{\CM}, \mathcal{C} \cup \{\textit{certified-edges}\}, \textit{pass})$.
The certified edges (those satisfying the claim) are committed.

\paragraph{Assert (failure).}
If $\textit{claim}$ evaluates to $\bot$:
$(\hat{\CM}, \mathcal{C}, \mathcal{R})
  \xrightarrow{\texttt{assert}(\textit{claim},\textit{on-fail}=f)}
  (\hat{\CM}, \mathcal{C}, \textit{narrow})[f]$,
where $[f]$ indicates that the fail-action $f$ is queued for execution.
The block does not halt; it narrows the scope and continues.

\paragraph{Recurse.}
$(\hat{\CM}, \mathcal{C}, \mathcal{R})
  \xrightarrow{\texttt{recurse}(\textit{cond}, \textit{body})}
  (\hat{\CM}^*, \mathcal{C}^*, \mathcal{R}^*)$,
where $(\hat{\CM}^*, \mathcal{C}^*, \mathcal{R}^*)$ is the fixed point of
repeatedly applying $\textit{body}$ until $\textit{cond}$ is satisfied.
Termination is guaranteed by the Cessation Theorem (Theorem~\ref{thm:cessation}).

\paragraph{Report.}
$(\hat{\CM}, \mathcal{C}, \mathcal{R})
  \xrightarrow{\texttt{report}(\textit{outputs})}
  (\hat{\CM}, \mathcal{C}, \mathcal{R})$
with side effect: emit the named outputs (contact map, regime classification,
inflation above floor) to the specified format.  The state is unchanged.
\end{definition}

\subsection{Semantic Properties}

\begin{theorem}[Compositional Soundness]
\label{thm:compositional}
Let $b_1, b_2$ be successive blocks in a DSL program with reduction
$\Gamma \xrightarrow{b_1} \Gamma' \xrightarrow{b_2} \Gamma''$.  Then:
\begin{enumerate}[label=(\roman*)]
  \item \textbf{Monotone commitment}: $\mathcal{C} \subseteq \mathcal{C}'
    \subseteq \mathcal{C}''$ (committed edges are never revoked).
  \item \textbf{Non-inflation}: $\hat{\CM}'' \leq \hat{\CM}' \leq \hat{\CM}$
    pointwise (estimates are monotone non-increasing in partition depth).
  \item \textbf{Floor preservation}: $\hat{\CM}'' \geq \CMbmin$ pointwise
    (no estimate can go below the floor).
\end{enumerate}
\end{theorem}

\begin{proof}
(i) Follows from the Assert rule: commitment extends $\mathcal{C}$ but never
removes from it.  The Non-Return Theorem (Theorem~\ref{thm:non-return})
applies.

(ii) Each Analyse and Compose step applies $\wedge$ (pointwise minimum),
so the estimate can only decrease or remain constant.

(iii) Every partition operation $\mathcal{O}$ produces estimates bounded
below by $\bmin$ (Theorem~\ref{thm:floor-positive}), so $\hat{\CM}_\mathcal{O}
\geq \CMbmin$.  The $\wedge$ of estimates bounded below by $\CMbmin$ is itself
bounded below by $\CMbmin$.
\end{proof}

\subsection{Execution Modes}

\begin{enumerate}[label=\textbf{Mode \arabic*.}, leftmargin=*]
\item \textbf{User-only.}  The user writes a complete DSL program including
  all \texttt{scope}, \texttt{analyse}, \texttt{compose}, \texttt{assert},
  and \texttt{report} blocks.  The CMG executes the program faithfully,
  performing the specified partition operations and returning the contact map
  as specified.  The system does not add any blocks.

\item \textbf{System-only.}  The user provides only a \texttt{scope} block
  (or a target path).  The CMG writes and executes the remainder of the
  program --- selecting methods, writing \texttt{analyse} and \texttt{compose}
  blocks, generating \texttt{assert} claims based on its current best estimate
  of $\CMbmin^\mathfrak{S}$, and recursing until the Cessation Theorem applies.

\item \textbf{Collaborative.}  The user writes \texttt{scope} and
  \texttt{assert} blocks (encoding architectural intuition and minimum
  individuation claims).  The CMG fills in the \texttt{analyse}, \texttt{compose},
  and \texttt{recurse} blocks.  Discrepancies between the user's assert claims
  and the CMG's computed contact map are surfaced as new \texttt{report} blocks
  indicating where the user's structural model diverges from the measured
  contact map.
\end{enumerate}

\subsection{Worked Examples}

\subsubsection{Example 1: Authentication Module Analysis}

The following script investigates the contact structure of an authentication
module.  It is a collaborative script: the user wrote \texttt{scope} and
\texttt{assert}; the CMG wrote \texttt{analyse}, \texttt{compose}, and the
\texttt{recurse} body.

\begin{lstlisting}[caption={Authentication module contact map investigation},
                   label={lst:auth}]
# Contact map investigation: auth module
# User-authored: scope, assert blocks
# System-authored: analyse, compose, recurse

scope authentication_module {
  target: "src/auth/**/*.ts"
  exclude: "**/*.test.ts"
  exclude: "**/*.d.ts"
  depth: 3
}

# System: extract structural partition operations
analyse static {
  extract: call_graph
  extract: type_apertures
  extract: dependency_edges
  extract: export_surface
}

# System: semantic inference of coordination boundaries
analyse semantic {
  method: hf
  model: "bigcode/starcoder2-7b"
  query: "identify coordination boundaries and
          interface contracts between components"
}

# System: compose static and semantic estimates
compose {
  merge: [static.call_graph,
          static.type_apertures,
          semantic.boundaries]
  resolve: overlaps by minimum_depth
}

# User: architectural individuation claim
# (TokenValidator and UserStore should be
#  at least 0.6 apart in S-entropy space)
assert individuation {
  claim: distance(auth.TokenValidator,
                  auth.UserStore) >= 0.6
  on_fail: narrow_scope(auth.TokenValidator,
                        auth.UserStore)
}

# User: regime claim (coherent or better)
assert regime {
  claim: R_est >= 0.8
  on_fail: recurse until marginal_return < floor {
    analyse dynamic {
      extract: runtime_call_trace
      extract: type_coercions
    }
    compose {
      merge: [current, dynamic.call_trace]
      resolve: overlaps by minimum_depth
    }
  }
}

# System: assert holonomy of all 3-cycles
assert individuation {
  claim: distance(auth.SessionManager,
                  auth.TokenValidator) >= 0.45
  on_fail: narrow_scope(auth.SessionManager,
                        auth.TokenValidator)
}

assert individuation {
  claim: distance(auth.UserStore,
                  auth.PasswordHasher) >= 0.55
  on_fail: narrow_scope(auth.UserStore,
                        auth.PasswordHasher)
}

report {
  output: contact_map
  output: regime_classification
  output: inflation_above_floor
  output: holonomy_per_cycle
  format: json
}
\end{lstlisting}

\subsubsection{Example 2: System-Only Investigation}

Here the user provides only a scope; the CMG generates the complete
investigation script.

\begin{lstlisting}[caption={System-generated contact map investigation},
                   label={lst:system-only}]
# Contact map investigation: payment service
# User-authored: scope only
# System-authored: all remaining blocks

scope payment_service {
  target: "services/payment/**"
  exclude: "**/vendor/**"
}

# System: layered partition analysis
analyse static {
  extract: call_graph
  extract: type_apertures
  extract: dependency_edges
  extract: api_surface
  extract: shared_state
}

analyse semantic {
  method: hf
  model: "bigcode/starcoder2-7b"
  query: "identify domain boundaries,
          transaction boundaries,
          and external service interfaces"
}

analyse static {
  method: holonomy
  extract: cycle_detection
  extract: holonomy_per_cycle
}

compose {
  merge: [static.call_graph,
          static.type_apertures,
          static.shared_state,
          semantic.domain_boundaries,
          static.holonomy_per_cycle]
  resolve: overlaps by minimum_depth
}

# System: floor-aware regime check
assert regime {
  claim: R_est >= 0.5
  on_fail: recurse until marginal_return < floor {
    analyse semantic {
      method: hf
      model: "bigcode/starcoder2-7b"
      query: "identify type coercions
              and implicit contracts"
    }
    compose {
      merge: [current,
              semantic.type_coercions]
      resolve: overlaps by minimum_depth
    }
  }
}

# System: assert minimum individuation for
# highest-tension edge pairs
assert individuation {
  claim: distance(payment.Processor,
                  payment.Ledger) >= floor
  on_fail: narrow_scope(payment.Processor,
                        payment.Ledger)
}

assert individuation {
  claim: distance(payment.Gateway,
                  payment.Validator) >= floor
  on_fail: narrow_scope(payment.Gateway,
                        payment.Validator)
}

report {
  output: contact_map
  output: regime_classification
  output: inflation_above_floor
  output: tension_per_edge
  output: holonomy_summary
  format: json
}
\end{lstlisting}

\subsubsection{Example 3: Discrepancy Surfacing}

This example illustrates a collaborative investigation where the user's
structural intuition (assert) is inconsistent with the measured contact map.
The CMG surfaces the discrepancy as a new report block.

\begin{lstlisting}[caption={Discrepancy between user model and measured contact map},
                   label={lst:discrepancy}]
# User believes the API layer and domain layer
# are well-separated (distance >= 0.8).
# CMG finds actual distance is 0.41.

scope api_domain_boundary {
  target: "src/**"
  exclude: "src/infrastructure/**"
}

analyse static {
  extract: call_graph
  extract: type_apertures
}

analyse semantic {
  method: hf
  model: "bigcode/starcoder2-7b"
  query: "characterise coupling between
          API handlers and domain objects"
}

compose {
  merge: [static.call_graph,
          semantic.coupling]
  resolve: overlaps by minimum_depth
}

# User assertion: architectural intent
assert individuation {
  claim: distance(api.handlers,
                  domain.aggregates) >= 0.8
  on_fail: narrow_scope(api.handlers,
                        domain.aggregates)
}

# CMG report: assertion failed.
# Measured distance: 0.41
# Regime: aperture-dominated (R_est = 0.38)
# The api.handlers layer directly references
# 14 domain.aggregates types, producing
# high S-entropy overlap.
# Inflation above floor: 1.93x

report {
  output: contact_map
  output: regime_classification
  output: inflation_above_floor
  output: discrepancy_report
  format: json
}
\end{lstlisting}

% ==========================================================================
\section{Discussion}
\label{sec:discussion}
% ==========================================================================

\subsection{Relation to Classical Testing}

Classical testing \citep{clarke1999model, baier2008principles} evaluates
predicates at points in a system's execution space.  Model checking
exhaustively verifies all paths through a finite state space; property-based
testing samples from a distribution of inputs.  Both paradigms are
\emph{point-evaluation} paradigms: the unit of information is a Boolean value
at a specific execution point.

The contact map framework is not a replacement for these paradigms within
their own domain of applicability.  It is a different level of analysis.
A test failure is \emph{evidence} that some pair of components is more
similar (in S-entropy space) than the minimum sufficient contact map would
predict --- that is, a test failure is a signal about the contact map, not
a refutation of a system-level property.  As established in
Section~\ref{sec:empirical}, every classical test is one instrument in the
array; no finite collection of such instruments determines $\CMbmin$.

More precisely: let $\phi(x, f(x)) = 0$ be a test failure at input $x$.
This event carries information $-\log_2 p_{\text{fail}}$ bits about the
execution.  In the contact map framework, this information is a partition
operation of type ``dynamic'' that updates the estimate $\hat{\CM}$.  The
test failure does not directly certify a partition depth; it contributes
evidence that the contact distance $\CM(u_i, u_j)$ between the failing unit
$u_i$ and the unit responsible for the violated property $u_j$ is less than
previously estimated.

\subsection{The First Cannot Answer the Second}

The central methodological claim of this paper is categorical:

\begin{proposition}
No finite collection of point-evaluation predicates can determine $\CMbmin$
for a non-trivial bounded resolvable space.
\end{proposition}

\begin{proof}
$\CMbmin$ is defined over the contact locus $\Part_{\bmin}$, which requires
comparison against all of $\Neg(A, \Part)$ for each component $A$.  By
Theorem~\ref{thm:incompletable}, no finite sequence of partition operations
completes this comparison.  A finite collection of point-evaluation
predicates is a finite sequence of partition operations of dynamic type.
Hence it cannot determine $\CMbmin$.  Moreover, by Rice's theorem
\citep{rice1953classes}, the semantic property defining $\CMbmin$ is
non-trivial and undecidable from local information alone.
\end{proof}

\subsection{The Second Subsumes the First}

\begin{proposition}
Every classical test assertion is a partition operation in the contact map
framework (specifically, a dynamic-type analyse block with boolean output).
\end{proposition}

\begin{proof}
A test assertion $\phi(x_i, f(x_i)) \in \{0,1\}$ can be encoded as a
partition operation $\mathcal{O}_{\phi,x_i} : \mathfrak{S} \to \hat{\CM}$
that updates the contact map estimate based on the pass/fail outcome.  A
passing test increases confidence that $\CM(u_i, u_j) \geq \theta$ for some
threshold $\theta$; a failing test provides evidence that $\CM(u_i, u_j) <
\theta$.  Hence every classical test is a special case of the partition
operations composable within the DSL's \texttt{analyse dynamic} block.
\end{proof}

\subsection{Falsifiability and Scope}

The framework is falsifiable in the following sense.

\begin{proposition}
The claim ``$\CMbmin^\mathfrak{S}(\cdot, \cdot) \geq \theta$ for all adjacent
pairs'' is falsified by exhibiting a pair $(u_i, u_j) \in E_C$ with
$\dS(S(u_i), S(u_j)) < \theta$, which is detectable by any partition
operation that discriminates $u_i$ from $u_j$ (under the discriminativity
condition on $\phi$).
\end{proposition}

The framework's scope is limited to systems that can be modelled as bounded
resolvable spaces --- i.e., finite software systems with well-defined unit
boundaries and measurable coupling weights.  It does not apply to systems
with zero-measure components (vacuous units) or systems with no measurable
separation (fully monolithic systems with no partition structure).

\subsection{The DSL as Audit Trail}

A completed DSL execution file serves as a complete audit trail for a contact
map investigation.  By the Non-Return Theorem (Theorem~\ref{thm:non-return})
and Compositional Soundness (Theorem~\ref{thm:compositional}), the sequence of
committed partitions in $\mathcal{C}$ is a monotone non-decreasing record of
the investigation's findings.  This has practical implications:

\begin{enumerate}[label=(\roman*)]
  \item \textbf{Reproducibility}: the script can be re-executed to reproduce
    the contact map estimate (modulo non-determinism in LLM inference);
  \item \textbf{Auditability}: the committed partitions form a verifiable
    record of what was individuated and when;
  \item \textbf{Extensibility}: new \texttt{analyse} blocks can be appended
    to refine the estimate without invalidating prior committed partitions.
\end{enumerate}

% ==========================================================================
\section{Conclusion}
\label{sec:conclusion}
% ==========================================================================

We have introduced a measure-theoretic framework for software system
characterisation grounded in the concept of the \emph{minimum sufficient
contact map}.  The framework rests on four main results.

First, the \emph{Partition Floor Theorem} (Theorem~\ref{thm:floor-positive}),
derived from the structure of individuation itself: an agent is an
individuated item; individuation by negation requires a separator of positive
measure; the separator is finite (because the agent is finite, by
Proposition~\ref{prop:agents-finite}); therefore $\bmin > 0$.  This
derivation supersedes earlier geometric proofs that contained unjustified
bounds.  The floor is not intrinsic to the space as an abstract object ---
it is the minimum cost of any single step of the incompletable individuation
comparison performed by any finite agent within the space.

Second, the \emph{Incompletability Theorem} (Theorem~\ref{thm:incompletable})
establishes that individuation of any item against its complement is never
finished by a finite agent in an infinite space.  This is the theoretical
core of the empirical observation (Section~\ref{sec:empirical}) that
scientific investigation never terminates: every measurement is a cell
subdivision, not a point approach.  Truth is a cell.

Third, the \emph{Minimum Sufficiency Trijection}
(Theorem~\ref{thm:minimum-sufficiency}) establishes a canonical bijective
correspondence among three objects: the partition floor, the minimum
sufficient contact map, and the minimum sufficient Lagrangian.  The minimum
sufficient action is a derived quantity (Corollary~\ref{cor:action}).  Every
additional structural annotation is a partition depth inflation; none reduces
the floor.

Fourth, the \emph{Local Invisibility Theorem}
(Theorem~\ref{thm:local-invisibility}) shows that universal local validity
does not imply global contact map coherence.  This is the formal basis for
the claim that classical point-evaluation testing cannot determine the
contact map --- a necessary limitation, not a contingent one.

The \emph{Wind Tunnel DSL}, specified by a formal context-free grammar with
compositional, monotone, floor-aware operational semantics, provides the
language in which contact map investigations are specified, executed, and
recorded.  Its co-authorship model --- users write structural intuition, the
system writes investigation --- is a practical instantiation of the
theoretical framework.

The \emph{Water Level Invariance Theorem} (Theorem~\ref{thm:water-beads})
formalises the core epistemic insight: rearranging observable states (test
outcomes, type annotations, coverage reports, instrument readings) does not
change the water level $\bmin$.  The task of the contact map generator is
to measure $\bmin$ accurately, not to remove it --- because removing it
would require the impossible: a finite agent completing an incompletable
comparison against the whole complement of the system.

Future work includes: (i) constructive bounds on $\bmin$ for specific
classes of software systems (microservices, monoliths, functional pipelines);
(ii) a type-theoretic embedding of the DSL in a dependent type system;
(iii) empirical calibration of the coordination regime boundaries against
known failure modes in production systems; (iv) extension to dynamically
evolving systems where the partition changes over time under the Non-Return
constraint.

% ==========================================================================
% Bibliography
% ==========================================================================

\bibliographystyle{plainnat}
\bibliography{references}

\end{document}
